\section{Toric filling of the unipotent cusp.}

At the unipotent puncture $p_0$, the local monodromy $T_0 = I + N$ has unipotent index $2$, with $N^2 = 0$ and $\operatorname{rk}(N) = 2$.
In this section we construct the compact analytic filling of the fibration over a punctured disc $\Delta^* = \{0 < |t_c| < r\}$ across $t_c = 0$, following Mumford's theory of toroidal degenerations of abelian varieties.

\subsection{The anticanonical hexagon of the degree-6 del Pezzo surface}

\begin{construction}\label{const:toric-model}
Let $N_{\dbR} \cong \dbR^2$ be the real vector space associated with the cocharacter lattice $\bar\Lambda = \Lambda/\Lambda_{\mathrm{tor}}$.
The fan $\Sigma_{A_2}$ of the complete degree-$6$ del Pezzo surface $dP_6 = \dbP^2 \sharp 3\overline{\dbP}^2$ consists of $6$ rays generated by the roots of the $A_2$ root system.
The anticanonical divisor $-K_{dP_6}$ is a cycle of $6$ smooth rational $(-1)$-curves:
\[
D = C_1 + C_2 + C_3 + C_4 + C_5 + C_6 \in |-K_{dP_6}|,
\]
meeting transversally in $6$ double points $q_1, \dots, q_6$.
The affine toric variety $U_{\Sigma} \to \dbC$ associated with the periodic cone decomposition over the lattice $\dbZ^2$ defines a family of abelian surfaces degenerating to the non-normal central fiber $W_0$.
\end{construction}

\begin{proposition}[Properties of the Central Fiber $W_0$]\label{prop:W0-properties}
The central fiber $W_0 = f^{-1}(p_0)$ constructed by the Mumford toroidal compactification satisfies:
\begin{enumerate}[label=(\roman*)]
    \item $W_0$ is a reduced, connected, non-normal complex analytic surface with normal crossings along the double curve locus $D_{\mathrm{sing}} = \mathrm{Sing}(W_0)$.
    \item The normalization $\nu \colon \widetilde W_0 \to W_0$ is the disjoint union of $6$ copies of $\dbP^1 \times \dbP^1$ or $dP_6$, glued along the anticanonical cycle $D$.
    \item The singular locus consists of $6$ rational curves intersecting at $6$ triple points.
    \item The topological Euler characteristic of $W_0$ is:
    \[
    e(W_0) = e(\widetilde W_0) - e(D_{\mathrm{sing}}) + e(\mathrm{Triple}) = 6(4) - 6(2) + 6 - 16 = 2.
    \]
\end{enumerate}
\end{proposition}

\begin{theorem}[Toric Collar Retraction]\label{thm:toric-retraction}
There exists a smooth collar neighborhood $N_0 \subset X$ containing $W_0$ such that:
\begin{enumerate}[label=(\alph*)]
    \item $N_0$ deformation retracts onto the singular fiber $W_0$: $r \colon N_0 \xrightarrow{\sim} W_0$.
    \item The boundary $\partial N_0$ is a Seifert fiber space over the circle $S^1$, diffeomorphic to the mapping torus of the unipotent monodromy $T_0 \in \mathrm{Sp}(4, \dbZ)$.
    \item The vanishing cycles in $H_2(F_b; \dbZ)$ generated by $\operatorname{im}(N) = \langle \gamma, u \rangle$ contract to $1$-cycles on $W_0$, matching the kernel of the specialization map.
\end{enumerate}
\end{theorem}
